% \documentclass{article}
% \usepackage[utf8]{inputenc}

% \title{TFG - DASIoT}
% \author{José Manuel Soto Hidalgo}
% \date{March 2022}

% \begin{document}

% \maketitle

% \section*{Objetivos}

% \section{Introduction}

% \section{Fundamentos}

% \section{Metodología}

% \subsection{Planificación a priori}

% \subsection{Planificación a posterior}

% \subsection{Presupuesto}

% \section{Conclusiones}

% \end{document}
\documentclass[spanish, 11pt, a4paper]{book}
\usepackage[T1]{fontenc} % Use 8-bit encoding that has 256 glyphs
\usepackage[utf8]{inputenc}
\usepackage{fourier} % Use the Adobe Utopia font for the document - comment this line to return to the LaTeX default
\usepackage{listings} % para insertar código con formato similar al editor
%\usepackage{babel} % Selecciona el español para palabras introducidas automáticamente, p.ej. "septiembre" en la fecha y especifica que se use la palabra Tabla en vez de Cuadro [spanish, es-tabla]
\usepackage[spanish,es-tabla]{babel}
\usepackage{url} % ,href} %para incluir URLs e hipervínculos dentro del texto (aunque hay que instalar href)
\usepackage{graphics,graphicx, float} %para incluir imágenes y colocarlas
\usepackage[gen]{eurosym} %para incluir el símbolo del euro
\usepackage{cite} %para incluir citas del archivo <nombre>.bib
\usepackage{enumerate}
\usepackage{hyperref}
\usepackage{graphicx}
\usepackage{tabularx}
\usepackage{booktabs}

\usepackage[table,xcdraw]{xcolor}
\hypersetup{
	colorlinks=true,	% false: boxed links; true: colored links
	linkcolor=black,	% color of internal links
	urlcolor=cyan		% color of external links
}
\renewcommand{\familydefault}{\sfdefault}
\usepackage{fancyhdr} % Custom headers and footers
\pagestyle{fancyplain} % Makes all pages in the document conform to the custom headers and footers
\fancyhead[L]{} % Empty left header
\fancyhead[C]{} % Empty center header
\fancyhead[R]{Luis Aróstegui Ruiz} % My name
\fancyfoot[L]{} % Empty left footer
\fancyfoot[C]{} % Empty center footer
\fancyfoot[R]{\thepage} % Page numbering for right footer
%\renewcommand{\headrulewidth}{0pt} % Remove header underlines
\renewcommand{\footrulewidth}{0pt} % Remove footer underlines
\setlength{\headheight}{13.6pt} % Customize the height of the header

\usepackage{titlesec, blindtext, color}
\definecolor{gray75}{gray}{0.75}
\newcommand{\hsp}{\hspace{20pt}}
\titleformat{\chapter}[hang]{\Huge\bfseries}{\thechapter\hsp\textcolor{gray75}{|}\hsp}{0pt}{\Huge\bfseries}
\setcounter{secnumdepth}{4}
\usepackage[Sonny]{fncychap}



\begin{document}

	% Plantilla portada UGR
	\begin{titlepage}
 
 
\newlength{\centeroffset}
\setlength{\centeroffset}{-0.5\oddsidemargin}
\addtolength{\centeroffset}{0.5\evensidemargin}
\thispagestyle{empty}

\noindent\hspace*{\centeroffset}\begin{minipage}{\textwidth}

\centering
\includegraphics[width=0.9\textwidth]{imagenes/logo_ugr.jpg}\\[1.4cm]

\textsc{ \Large TRABAJO FIN DE GRADO\\[0.2cm]}
\textsc{ INGENIERÍA EN ...}\\[1cm]
% Upper part of the page
% 
% Title
{\Huge\bfseries Titulo del Proyecto\\
}
\noindent\rule[-1ex]{\textwidth}{3pt}\\[3.5ex]
{\large\bfseries Subtitulo del Proyecto}
\end{minipage}

\vspace{2.5cm}
\noindent\hspace*{\centeroffset}\begin{minipage}{\textwidth}
\centering

\textbf{Autor}\\ {Nombre Apellido1 Apellido2 (alumno)}\\[2.5ex]
\textbf{Directores}\\
{Nombre Apellido1 Apellido2 (tutor1)\\
Nombre Apellido1 Apellido2 (tutor2)}\\[2cm]
\includegraphics[width=0.3\textwidth]{imagenes/etsiit_logo.png}\\[0.1cm]
\textsc{Escuela Técnica Superior de Ingenierías Informática y de Telecomunicación}\\
\textsc{---}\\
Granada, mes de 201
\end{minipage}
%\addtolength{\textwidth}{\centeroffset}
%\vspace{\stretch{2}}
\end{titlepage}



	
	% Plantilla prefacio UGR
	\thispagestyle{empty}

\begin{center}
{\large\bfseries TFG - DASIoT: Desarrollo y Auditoría de Seguridad para prototipo de dispositivos IoT}\\
\end{center}
\begin{center}
Luis Aróstegui Ruiz\\
\end{center}

%\vspace{0.7cm}
\vspace{0.5cm}
\noindent{\textbf{Palabras clave}: Internet de las Cosas, Seguridad}\\

\vspace{0.7cm}
\noindent{\textbf{Resumen}}\\

En el contexto de IoT, hay un ecosistema de entornos de desarrollo específicos. En este TFG se hará uso de alguno de ellos para llevar a cabo la implementación de un dispositivo IoT y analizar los potenciales problemas de seguridad que pueden aparecer durante la etapa de implementación.

% El Internet de las Cosas (IoT) se utiliza para proporcionar conectividad entre diferentes
% dispositivos, los cuales están conectados a Internet. Es un sistema donde los objetos
% actúan con otros objetos a través de un medio de comunicación inalámbrico para
% intercambiar y transferir información sin interacción humana.
% Los dispositivos IoT son propensos a ataques vulnerables debido a la naturaleza simple
% y abierta de sus redes. Por lo tanto, la privacidad y la seguridad son la mayor
% preocupación de esta tecnología. El enfoque de las amenazas de seguridad y privacidad
% en IoT es crucial para promover el desarrollo de IoT.
% Este trabajo final de grado tiene como objetivo conceptualizar el Internet de las Cosas,
% establecer cuáles son sus principales características, y entrar en detalle en los aspectos
% de la privacidad y la seguridad de los datos personales. Cuáles son esas amenazas, la
% problemática particular del Internet de las Cosas que las genera y como mitigar el riesgo
% a dichas amenazas.
% En este trabajo nos haremos las siguientes preguntas: ¿Qué pasa con nuestros datos
% personales? ¿A qué riesgos estamos expuestos aportando tanta información “privada”?
% ¿Qué seguridad existe en el Internet de las Cosas? En este proyecto se pretende
% responder a estas y a muchas otras preguntas. Para ello vamos a profundizar en el tema
% de la privacidad y seguridad en el Internet de las Cosas.

\cleardoublepage


\begin{center}
{\large\bfseries Project Title: Project Subtitle}\\
\end{center}
\begin{center}
First name, Family name (student)\\
\end{center}

%\vspace{0.7cm}
\noindent{\textbf{Keywords}: Keyword1, Keyword2, Keyword3, ....}\\

\vspace{0.7cm}
\noindent{\textbf{Abstract}}\\

Write here the abstract in English.

\cleardoublepage
\thispagestyle{empty}

\noindent\rule[-1ex]{\textwidth}{2pt}\\[4.5ex]

Yo, \textbf{Nombre Apellido1 Apellido2}, alumno de la titulación TITULACIÓN de la \textbf{Escuela Técnica Superior
de Ingenierías Informática y de Telecomunicación de la Universidad de Granada}, con DNI XXXXXXXXX, autorizo la
ubicación de la siguiente copia de mi Trabajo Fin de Grado en la biblioteca del centro para que pueda ser
consultada por las personas que lo deseen.

\vspace{6cm}

\noindent Fdo: Nombre Apellido1 Apellido2

\vspace{2cm}

\begin{flushright}
Granada a X de mes de 201 .
\end{flushright}

\cleardoublepage
\thispagestyle{empty}

\noindent\rule[-1ex]{\textwidth}{2pt}\\[4.5ex]

D. \textbf{Nombre Apellido1 Apellido2 (tutor1)}, Profesor del Área de XXXX del Departamento YYYY de la Universidad de Granada.

\vspace{0.5cm}

D. \textbf{Nombre Apellido1 Apellido2 (tutor2)}, Profesor del Área de XXXX del Departamento YYYY de la Universidad de Granada.


\vspace{0.5cm}

\textbf{Informan:}

\vspace{0.5cm}

Que el presente trabajo, titulado \textit{\textbf{Título del proyecto, Subtítulo del proyecto}},
ha sido realizado bajo su supervisión por \textbf{Nombre Apellido1 Apellido2 (alumno)}, y autorizamos la defensa de dicho trabajo ante el tribunal
que corresponda.

\vspace{0.5cm}

Y para que conste, expiden y firman el presente informe en Granada a X de mes de 201 .

\vspace{1cm}

\textbf{Los directores:}

\vspace{5cm}

\noindent \textbf{Nombre Apellido1 Apellido2 (tutor1) \ \ \ \ \ Nombre Apellido1 Apellido2 (tutor2)}

\chapter*{Agradecimientos}
\thispagestyle{empty}

       \vspace{1cm}


Poner aquí agradecimientos...


	
	% Índice de contenidos
	\newpage
	\tableofcontents

	% Índice de imágenes y tablas
	\newpage
	\listoffigures

	% Si hay suficientes se incluirá dicho índice
	\listoftables 
	\newpage
	

	% Introducción 
	\chapter{Introducción}
	
	% Planificación
	\chapter{Estado del arte}

{\color{blue}

Antes de empezar a desarrollar los objetivos del proyecto es necesario conocer primero los fundamentos del Internet de las Cosas, esto engloba desde su funcionamiento, tipos de arquitecturas hasta los distintos estándares que existen. % Añadir parte de seguridad (normativas)

\section{Internet de las Cosas}

El Internet de las Cosas %(Ngu, Gutierrez, Metsis, Nepal, y Sheng, 2016)
permite que las personas
interactúen con dispositivos de distinto tipo conectados a Internet, como sensores o actuadores. Con
esto, se consigue integrar cualquier objeto del mundo real a Internet, estableciendo relaciones entre
todos ellos y cambiando la forma en la que interaccionamos con nuestro entorno. \\

La premisa básica y el objetivo de IoT es "conectar lo que no está conectado". Esto significa que los objetos que no están actualmente unidos a una red informática, es decir, a Internet, se conectarán para que puedan comunicarse e interactuar con personas y otros objetos. IoT es una transición tecnológica en la que los dispositivos nos permitirán sentir y controlar el mundo físico haciendo que los objetos sean más inteligentes y conectándolos a través de una red inteligente. 1 Cuando los objetos y las máquinas pueden ser detectados y controlados a distancia a través de una red, se consigue una mayor integración entre el mundo físico y los ordenadores. Esto permite mejoras en las áreas de eficiencia, precisión, automatización y habilitación de aplicaciones avanzadas. \\

El mundo del IoT es amplio y puede resultar algo complicado al principio debido a la abundancia de componentes y protocolos que engloba. En lugar de considerar IoT como un único tecnología, es bueno verlo como un paraguas de varios conceptos, protocolos y tecnologías, todos ellos dependientes a veces de un sector concreto. Aunque la amplia gama de elementos de IoT está diseñado para crear numerosos beneficios en las áreas de productividad y automatización, al mismo al mismo tiempo, introduce nuevos retos, como la ampliación del gran número de dispositivos y de la cantidad de datos que deben procesarse. que hay que procesar. \cite{hanes2017iot}

}

	% Planificación
	\chapter{Planificación}

\section{Planificación a priori}

	% Análisis del problema
	\chapter{Análisis del problema}

	% Desarrollo
	\chapter{Diseño}
	
	% Implementación
	\chapter{Implementación}

	% Pruebas
	\chapter{Pruebas}

	% Conclusiones
	\chapter{Conclusiones y trabajos futuros}


	
	\newpage
	\bibliography{bibliografia}
	\bibliographystyle{plain}
	
\end{document}